\documentclass[a4paper,11pt]{article}

\usepackage[a4paper, top=2.5cm, bottom=2.5cm]{geometry}
\usepackage[utf8]{inputenc}
\usepackage[french, english]{babel}

\selectlanguage{french}

\usepackage{enumitem}
\usepackage{hyperref}
\usepackage{multicol}
\usepackage{float}
\usepackage{array}
\usepackage{caption}
\usepackage{listings}

\renewcommand{\thesection}{\arabic{section}\quad---}
\renewcommand{\thesubsection}{\arabic{section}.\arabic{subsection}\quad---}

\newcolumntype{M}[1]{>{\raggedright}m{#1}}
\newcolumntype{C}[1]{>{\centering\arraybackslash }m{#1}}

%BEGIN LATEX
  \usepackage{core}

  \newcommand{\mousevox}[2]{\IfStrEq{\jobname}{\detokenize{mouse}}{#1}{#2}}
%END LATEX

\title{
  \blogtitle{Du Multi-Seat Linux}
}
\author{}
\date{}

\pagestyle{plain}
\begin{document}
\neobar{article06.fr}{nav/navi.fr}

\maketitle

\thispagestyle{first}
\begin{blogmain}
Lorsque vous étiez adolescents, peut-être avez-vous déjà dû coordonner des horaires d'utilisation d'une console de jeux ou d'une station de travail.
C'est une expérience frustrante d'attendre la manette ou le clavier, c'est une situation qui devrait pouvoir être évitée.
Bien-sûr, certains jeux coopératifs comme \href{https://supertuxkart.net/Main_Page.html}{SuperTuxKart} proposent du multi-manette et de l’écran partagé via du split-screen, cependant d'autres comme \href{https://www.luanti.org/fr/}{Luanti} n'ont pas cette fonctionnalité de split-screen.
Certains gestionnaires de fenêtre Linux comme Sway intègrent du multi-curseur à focus indépendants, cependant l'expérience reste très exotique et ne va encore plus loin comme du multi-claviers à focus indépendants.
\\
La solution que je vais vous proposer est celle du multi-seat par utilisateur.

\simpleschemas{1.0}{multiseat}{Multi-seat de 2 utilisateurs jouant à Luanti sur la même station de travail.}

Le projet Freedesktop fondé pour l'interopérabilité des environnements de bureau sous Unix défini que :
\begin{itemize}
  \item Le \href{https://www.freedesktop.org/wiki/Software/systemd/multiseat#definitionofterms}{multi-seat} est un système permettant de définir plusieurs ensembles indépendants de seats utilisables simultanément.
  \item Le concept de \href{https://www.freedesktop.org/wiki/Software/systemd/multiseat#definitionofterms}{seat} est un groupe de matériels comprenant au moins une interface graphique, comprenant généralement un clavier et une souris.
\end{itemize}

Sous Linux, le multi-seat est communément utilisé par des gestionnaires de sessions comme LightDM pour attribuer des sièges à des sessions.
Cela c'est plus répandu dans le début des années 2010 avec la généralisation de systemd et de sa brique systemd-login qui a formalisé le concept de seat.

Aujourd'hui, avec les progrès de bibliothèques graphiques comme la SDL3 et sa récente \href{https://github.com/libsdl-org/SDL/pull/12626}{intégration du support Multi-Seat} documenté par l'article \href{https://www.phoronix.com/news/SDL-Merges-Wayland-Multi-Seat}{SDL Merges Wayland Multi-Seat Support}, nous pouvons plus facilement programmer des applications multi-utilisateurs simultanées.
Ces applications peuvent par exemple être des jeux coopératifs, et des applications de dessins coopératifs au futur.

\section{Du multi-seats à la Nix}
Le scénario d'illustration que je vous propose est celui où nous définirons un seat par utilisateur de Sway.

L'article \href{https://wiki.nixos.org/wiki/Multiseat}{Multiseat} du Wiki de NixOS nous propose de définir des règles udev pour attribuer des interfaces à nos seats, et nous mentionne LightDM comme gestionnaire de session.

Nous pouvons définir nos seats et leurs interfaces comme suit :
\includesnippet{udev-rules}{Règles udev de 2 seats}
% le seat0 aura un iGPU et un clavier/des souris, le seat1 aura un GPU et un clavier/une souris + interface audio.
% loginctl seat-status seat0 seat1

\newpage
Nous pouvons configurer LightDM pour attribuer des seats aux utilisateurs comme suit :
\includesnippet{lightdm}{LightDM de 2 utilisateurs}

\section{Retour d'expérience}
J'ai trouvé que ce sujet était peu documenté sous Linux.
J'utilise du multi-seat pour proposer à mes amis de passage un second siège avec une session invitée.
Nous pouvons partager la même distribution Linux ou bien ceux-ci peuvent chrooter dans la leurs.

Ce modèle est pas mal du tout et est agréable à utiliser. Nous pouvons y ajouter des services de partage de curseurs et claviers comme \href{https://github.com/feschber/lan-mouse/tree/main}{Lan Mouse} pour rendre l'expérience plus collaborative.
\end{blogmain}
\end{document}
